\makeabstract
    {
    Was Plate Tectonics Active During the Archean Eon? A Paleomagnetic Approach
    }
    {
    Drew Weisgerber\textsuperscript{1,2}, Peter Blatchford\textsuperscript{2}, Roger Fu\textsuperscript{2}
    }
    {
    \textsuperscript{1} Amherst College, Amherst, MA\\
\textsuperscript{2} Department of Earth and Planetary Sciences, Harvard University, Cambridge, MA
    }
    {
    The theory of plate tectonics is the central paradigm of modern geology, as it provides a framework for understanding the history and morphology of continental and oceanic crust, as well as the distribution of volcanoes and seismic zones on the modern Earth. While there is evidence that plate tectonics has been active throughout the Phanerozoic, there is no consensus about when it began, nor what characterized the lithospheric regime that preceded it. Determining the nature and timing of the onset of plate tectonics is essential to understanding ancient climate and the development of life, as tectonic processes play an essential role in the modern carbon cycle.
    }
    {
    This study investigates ancient tectonics through paleomagnetic analysis of $\sim$3.1 Ga volcanic rocks of the Sholl Terrane, collected from the Northwest Pilbara Craton in Australia. The orientation of the ancient magnetic field was determined by measuring the magnetic moment of oriented rock samples in a 2G Enterprises DC-SQuID Superconducting Rock Magnetometer. Through progressive thermal demagnetization past the blocking temperatures of the dominant ferromagnetic mineral phases, the characteristic remanent magnetism (ChRM) component was isolated.
    }
    {
    ChRM orientations will be used to define a paleomagnetic pole for the Northwest Pilbara Craton at $\sim$3.1 Ga. This will be compared to paleomagnetic poles from adjacent units to quantify the extent of apparent polar wander (APW) and evaluate the possibility of mobile-lid tectonics in the Archean.
    }
    {
    Determining the nature and timing of the onset of plate tectonics is essential to understanding ancient climate and the development of life, as tectonic processes play an essential role in the modern carbon cycle. This research seeks to shed light on the early tectonic history of the Earth, helping to explain unanswered questions about our planet's physical and biological history.
    }

    \newpage
    